\documentclass[11pt]{article}
\usepackage[a4paper,margin=26mm]{geometry}
\usepackage{amsmath,amssymb}
\setlength{\parindent}{0pt}
\setlength{\parskip}{7pt}
\emergencystretch=2em
\title{A metric spectral midpoint is not an optimal gradient step\\
\large Disproof of conjecture 00000008854}
\author{ziangni-sys}
\date{4 October 2026}
\begin{document}
\maketitle
The English conjecture asserts that the optimal step size is the midpoint
of the spectral interval of the metric matrix. The Chinese version likewise
specifies the metric matrix. We use the standard metric projected-gradient
interpretation: for a positive definite metric $H$ and constraint set $C$,
\[
 x_{k+1}=P_C^H\bigl(x_k-\eta H^{-1}\nabla f(x_k)\bigr).
\]
Here $P_C^H(y)$ minimizes $(z-y)^\top H(z-y)$ over $z\in C$.
The example is one dimensional, has a unique attained minimizer, and
uses a smooth strongly convex objective. Its metric is the identity, so
inverse-metric convention choices make no difference.

\textbf{Actual metric and projection.}
Take the real Hilbert space $\mathbb R$, $C=\mathbb R$, and $H=[1]$.
This is symmetric positive definite. The eigenvalue equation $Hv=\lambda v$
with $v\ne0$ gives $\lambda=1$, and $v=1$ witnesses this eigenvalue.
Hence its actual spectrum is $\{1\}$, its spectral interval is $[1,1]$,
and the proposed midpoint is
\[
 \eta_{\rm mid}=(1+1)/2=1.
\]
The metric projection is exactly $P_C^H(y)=y$: the objective
$(z-y)^2$ is nonnegative, equals zero at $z=y$, and vanishes only there.

\textbf{Objective, gradient, and iteration.}
Let $f(x)=2x^2$. Its actual derivative and Hilbert gradient are $4x$.
Its second derivative is $4>0$; directly, for every $x,y$,
\[
 f(y)=f(x)+4x(y-x)+2(y-x)^2.
\]
This identity gives strong convexity with curvature $4$.
The gradient is globally $4$-Lipschitz and the unique minimizer is zero.
Since the inverse metric and projection are identities, the actual update is
\[
 G_\eta(x)=x-\eta(4x)=(1-4\eta)x.
\]
For initial value $x_0=1$, induction yields the complete orbit
\[
 x_n=(1-4\eta)^n.
\]

\newpage
\textbf{The proposed step diverges; the optimal step terminates.}
For the spectral midpoint $\eta=1$, the orbit is $x_n=(-3)^n$.
Its distance from the unique minimizer is
\[
 |x_n|=3^n\longrightarrow+\infty.
\]
Thus the proposed midpoint is not even a convergent step in this
smooth strongly convex example.

In contrast, $\eta=1/4$ gives $G_\eta(x)=0$ for every initial value,
and every iterate after the first is the exact minimizer. More precisely,
the global one-step error factor is
\[
 q(\eta)=|1-4\eta|.
\]
It is nonnegative and has its unique minimum zero at $\eta=1/4$.
Indeed $q(\eta)=0$ if and only if $1-4\eta=0$. For every $x$,
$|G_\eta(x)|=q(\eta)|x|$, so this is the actual error contraction
factor, not an unrelated scalar surrogate.
The midpoint factor is $q(1)=3$, whereas $q(1/4)=0$.
Even among nonnegative step sizes, the midpoint cannot be optimal.

This disproves the metric-matrix midpoint clause. The example distinguishes
the metric matrix from the objective's Hessian: their spectra are
$\{1\}$ and $\{4\}$ respectively. A valid step-size rule must account for
the objective's curvature relative to the metric. We do not separately
refute the other convergence-rate clause or assert a theorem for an
unspecified meaning of generalized projection.

\textbf{Formal correspondence.}
The Lean project uses the actual identity continuous linear map on
$\mathbb R$ as the metric and the actual one-by-one identity matrix.
\texttt{matrix\_representation} links their actions.
\texttt{actual\_spectrum} characterizes the full set of eigenvalues by
nonzero eigenvectors; in finite dimension this is the metric spectrum.
\texttt{actual\_projection} and \texttt{projection\_unique} verify the
global metric minimization definition.

\texttt{actual\_derivative} and \texttt{actual\_gradient} use Mathlib's
derivative and gradient predicates. \texttt{unique\_minimum} proves
global minimizer uniqueness. \texttt{step\_formula} and
\texttt{actual\_iterates} prove the update and every iterate.
\texttt{midpoint\_escape} proves divergence of the actual absolute
errors to infinity. \texttt{quarter\_termination} proves exact termination,
and \texttt{unique\_optimal\_step} proves global optimality of the
error factor. The final \texttt{counterexample} assembles these objects.

\textbf{Reproduction.}
Use Lean 4.19.0 with public Mathlib revision
\begin{center}\small\texttt{c44e0c8ee63ca166450922a373c7409c5d26b00b}\end{center}
Run \texttt{lake build} inside \texttt{lean/}. The project prints
axiom audits; the validation notes record the completed build and
compiled PDF inspection.
\end{document}
